\begin{samepage}
\section{Conclusion}
We propose IntALNS for cooperative marine monitoring, formulating task allocation and route sequencing under platform capability, service-start time-window, and transfer-energy constraints. We design an agent-control mechanism that integrates allocation observations, execution history, and strategy examples to guide feasible ALNS route reconstruction. Requiring rationales to cite actual historical feedback makes search adjustments interpretable, as illustrated by the recorded case. Across \PairCount{} synthetic conditions, IntALNS reduces mean priority loss by \PriorityReduction\% and increases complete allocations from \BasicComplete{} to \AgentComplete{} compared with Basic under the same search budget. IntALNS thus combines better allocation of important monitoring tasks with search control grounded in recorded evidence.
\par\end{samepage}
